\chapter{Multi-Armed Bandits}
\label{ch:bandits}

\begin{quote}
\itshape
The simplest setting in which the exploration--exploitation dilemma already appears in full force is the multi-armed bandit problem.
\end{quote}

\bigskip

In this chapter we study the \emph{multi-armed bandit}\index{multi-armed bandit} problem\,---\,a sequential decision-making task without states, where the agent repeatedly chooses one of $k$ actions and observes a stochastic reward.  Despite its simplicity, the bandit setting already exposes the core tension between gathering information (exploration) and maximising immediate payoff (exploitation).  We begin with the basic formulation, analyse $\varepsilon$-greedy strategies, and then develop more principled approaches: \emph{Upper Confidence Bounds} (UCB), \emph{Thompson Sampling}, and the information-theoretic lower bound of Lai and Robbins.  The chapter closes with a discussion of contextual bandits\,---\,a bridge toward the full reinforcement learning problem\,---\,and an empirical comparison of the algorithms studied.


% ===========================================================
\section{Problem Formulation}
\label{sec:bandit-formulation}

\begin{definition}[$k$-armed bandit]
\label{def:k-armed-bandit}
\index{$k$-armed bandit}
A \textbf{$k$-armed bandit} is a sequential decision problem in which, at each time step $t=1,2,\dots$, the agent selects an action $A_t$ from a finite set
\[
\cA = \{1,2,\dots,k\},
\]
and receives a stochastic reward $R_t\in\R$.
\end{definition}

\begin{remark}[Notation convention]
In the bandit setting, we write $R_t$ for the reward received at time~$t$ in response to action~$A_t$.  This differs from the MDP convention introduced in Chapter~\ref{ch:mdp}, where the reward is denoted~$R_{t+1}$ (received after transitioning from state~$S_t$ with action~$A_t$).  The distinction arises because bandits have no state transitions; the simpler indexing $R_t$ is standard in the bandit literature.
\end{remark}

For each action $a\in\cA$ there is a reward distribution $\nu_a$.  When action $A_t=a$ is chosen, the reward is drawn independently from $\nu_a$.  In the \emph{stationary} setting these distributions do not change over time.

Unlike the general RL problem studied later in this book, the bandit setting has:
\begin{itemize}[leftmargin=*]
\item no notion of state;
\item immediate (non-delayed) rewards;
\item no transition dynamics\,---\,each round is independent.
\end{itemize}

Nevertheless, the fundamental challenge is already present: to learn which action is best, the agent must try different actions, yet to maximise reward it should favour actions that already seem good.


% ===========================================================
\section{Action Values, Optimal Actions, and Regret}
\label{sec:action-values-regret}

\begin{definition}[Action value]
\label{def:action-value-bandit}
\index{action value}
The \textbf{true value} (or \emph{expected reward}) of action $a\in\cA$ is
\begin{equation}
\label{eq:true-value}
q_*(a) = \E[R_t \mid A_t = a].
\end{equation}
\end{definition}

If the values~\eqref{eq:true-value} were known, the optimal strategy would simply be to always play
\begin{equation}
\label{eq:optimal-action}
a^* \in \argmax_{a\in\cA}\; q_*(a).
\end{equation}
Let
\begin{equation}
\label{eq:optimal-value}
q^* = \max_{a\in\cA}\; q_*(a)
\end{equation}
denote the value of the best action.

Since the values are unknown, we measure performance by comparison with the omniscient strategy.

\begin{definition}[Pseudo-regret]
\label{def:pseudo-regret}
\index{pseudo-regret}\index{regret}
The \textbf{pseudo-regret} after $T$ steps is
\begin{equation}
\label{eq:regret}
\cR_T = T\,q^* - \sum_{t=1}^{T} q_*(A_t).
\end{equation}
Its expectation equals
\[
\E[\cR_T] = T\,q^* - \sum_{t=1}^{T}\E[q_*(A_t)].
\]
\end{definition}

Intuitively, $\cR_T$ measures how much reward the algorithm ``left on the table'' compared to always playing the optimal arm.

Define the \emph{sub-optimality gap}\index{sub-optimality gap} of action $a$ as
\begin{equation}
\label{eq:gap}
\Delta_a = q^* - q_*(a) \ge 0,
\end{equation}
and let $N_T(a) = \sum_{t=1}^{T}\ind\{A_t = a\}$ count the number of times action $a$ has been played.

\begin{proposition}[Regret decomposition]
\label{prop:regret-decomp}
\index{regret decomposition}
In the stationary bandit setting (where the reward distributions do not change over time),
\begin{equation}
\label{eq:regret-decomposition}
\E[\cR_T] = \sum_{a\in\cA} \Delta_a\,\E[N_T(a)].
\end{equation}
\end{proposition}

\begin{proof}
Because the setting is stationary, $q_*(a)$ is constant across rounds.
Since $\sum_{a\in\cA}N_T(a)=T$ and $\sum_{t=1}^{T}q_*(A_t)=\sum_{a\in\cA}q_*(a)\,N_T(a)$ (grouping the sum by which arm was played), we have
\[
\cR_T
= T\,q^* - \sum_{a\in\cA}q_*(a)\,N_T(a)
= \sum_{a\in\cA}(q^*-q_*(a))\,N_T(a)
= \sum_{a\in\cA}\Delta_a\,N_T(a).
\]
Taking expectations on both sides yields~\eqref{eq:regret-decomposition}.
\end{proof}

Formula~\eqref{eq:regret-decomposition} is fundamental: minimising regret is equivalent to limiting the number of times sub-optimal actions are played.

\begin{remark}
If $\E[\cR_T]=O(\log T)$, then the average per-step regret decays as $O(\log T/T)\to 0$, meaning the agent chooses sub-optimal actions with vanishing frequency.
\end{remark}


% ===========================================================
\section{Greedy and \texorpdfstring{$\varepsilon$}{epsilon}-Greedy Strategies}
\label{sec:eps-greedy}

Since true action values are unknown, the agent must estimate them from data.

\begin{definition}[Action-value estimate]
\label{def:action-value-estimate}
\index{action-value estimate}
The \textbf{action-value estimate} $Q_t(a)$ is a quantity computed from all rewards observed before step $t$ for action $a$.
\end{definition}

\begin{definition}[Greedy action]
\label{def:greedy-action}
\index{greedy action}
An action $a$ is \textbf{greedy} at step $t$ if
\begin{equation}
\label{eq:greedy-action}
a \in \argmax_{b\in\cA}\; Q_t(b).
\end{equation}
\end{definition}

The simplest strategy is \emph{pure greedy}\index{pure greedy}: always choose the greedy action.  While intuitive, this can fail catastrophically: early noisy observations may cause the agent to underestimate the best arm and never return to it.

\begin{example}
\label{ex:greedy-stuck}
Consider a $2$-armed bandit with $q_*(1)=1$ and $q_*(2)=0.9$.  If the first observation from arm~1 happens to be unusually low and the first from arm~2 is unusually high, a greedy agent may lock onto arm~2 and never re-explore arm~1.
\end{example}

To escape such traps, we introduce random exploration.

\begin{definition}[$\varepsilon$-greedy strategy]
\label{def:eps-greedy}
\index{$\varepsilon$-greedy}
The \textbf{$\varepsilon$-greedy strategy} selects actions by the rule
\begin{equation}
\label{eq:eps-greedy}
A_t =
\begin{cases}
\displaystyle\argmax_{a\in\cA}Q_t(a), & \text{with probability } 1-\varepsilon,\\[6pt]
\text{uniform random from }\cA, & \text{with probability } \varepsilon,
\end{cases}
\end{equation}
where $\varepsilon\in[0,1]$.
\end{definition}

\begin{remark}
\label{rem:eps-greedy-regret}
If the optimal action is correctly identified as greedy, then it is selected with probability at least $1-\varepsilon+\varepsilon/k$.  However, with a fixed $\varepsilon>0$, the expected number of random (potentially sub-optimal) choices grows linearly in $T$, yielding $\E[\cR_T]=\Omega(T)$.  That is, fixed-$\varepsilon$ strategies have \emph{linear} regret.
\end{remark}


% ===========================================================
\section{Sample Averages and Incremental Updates}
\label{sec:incremental}

The most natural estimate of $q_*(a)$ is the \emph{sample average}\index{sample average}:
\begin{equation}
\label{eq:sample-average}
Q_t(a) = \frac{\sum_{i=1}^{t-1}R_i\,\ind\{A_i=a\}}{N_t(a)}, \qquad N_t(a)>0.
\end{equation}
For actions not yet tried, we initialise $Q_1(a)=0$.

\begin{proposition}
If the rewards for action $a$ are i.i.d.\ with mean $q_*(a)$ and $N_T(a)\to\infty$, then $Q_T(a)\to q_*(a)$ almost surely, by the strong law of large numbers.
\end{proposition}

Computing~\eqref{eq:sample-average} from scratch at every step is wasteful.  Instead, we use the \emph{incremental update}\index{incremental update}:
\begin{equation}
\label{eq:incremental-update}
Q_{n+1}(a) = Q_n(a) + \frac{1}{n+1}\bigl(R_{n+1}(a) - Q_n(a)\bigr),
\end{equation}
where $n=N_t(a)$ and $R_{n+1}(a)$ is the newest reward for action $a$.

\begin{remark}[General update rule]
\label{rem:general-update}
\index{step size}
Equation~\eqref{eq:incremental-update} is a special case of
\begin{equation}
\label{eq:general-update}
Q \leftarrow Q + \alpha\,(R - Q),
\end{equation}
with $\alpha=1/(n+1)$.  For \emph{non-stationary} problems a constant step size $\alpha\in(0,1)$ is often preferred, giving exponentially decaying weights to older observations:
\[
Q_{n+1} = (1-\alpha)^n Q_1 + \sum_{i=1}^{n}\alpha(1-\alpha)^{n-i}R_i.
\]
\end{remark}

The $\varepsilon$-greedy bandit agent can now be stated concisely:

\begin{algorithm}[ht]
\caption{$\varepsilon$-Greedy Bandit Agent}
\label{alg:eps-greedy}
\KwIn{Number of arms $k$, exploration rate $\varepsilon$, horizon $T$}
Initialise $Q(a)\leftarrow 0$, $N(a)\leftarrow 0$ for all $a\in\cA$\;
\For{$t=1,2,\dots,T$}{
  With probability $1-\varepsilon$: $A_t\leftarrow\argmax_{a}Q(a)$\;
  With probability $\varepsilon$: $A_t\leftarrow\text{Uniform}(\cA)$\;
  Observe reward $R_t$\;
  $N(A_t)\leftarrow N(A_t)+1$\;
  $Q(A_t)\leftarrow Q(A_t)+\frac{1}{N(A_t)}(R_t - Q(A_t))$\;
}
\end{algorithm}


% ===========================================================
\section{Upper Confidence Bounds (UCB)}
\label{sec:ucb}
\index{UCB}\index{Upper Confidence Bound}

The $\varepsilon$-greedy strategy explores \emph{uniformly at random}, without regard for which arms are more uncertain.  A more principled approach is to explore preferentially toward actions about which we know the least.  This is the idea behind \emph{optimism in the face of uncertainty}\index{optimism in the face of uncertainty}.

\subsection*{The Principle of Optimism}

The intuition behind optimism is elegant: when we are uncertain about the value of an arm, we should \emph{assume the best}.  If that optimistic assumption turns out to be correct, we are rewarded handsomely.  If it turns out to be wrong\,---\,the arm is not as good as we hoped\,---\,then we have gained useful information and will assign that arm a lower confidence bound in the future.  Either way, optimism is productive.

More concretely, suppose we maintain for each arm $a$ an upper confidence bound $U_t(a)$ such that $q_*(a) \le U_t(a)$ holds with high probability.  Selecting the arm with the largest $U_t(a)$ simultaneously exploits arms that are genuinely good and explores arms whose true value is still uncertain enough that their upper bound is elevated.  As an arm is pulled more, observations accumulate and the bound $U_t(a)$ tightens around the empirical mean $Q_t(a)$; an arm that turns out to be sub-optimal will eventually have a low upper bound and be abandoned.

\begin{definition}[UCB1]
\label{def:ucb1}
\index{UCB1}
The \textbf{UCB1} algorithm (Auer, Cesa-Bianchi, and Fischer, 2002) selects the action
\begin{equation}
\label{eq:ucb1}
A_t = \argmax_{a\in\cA}\left[Q_t(a) + c\,\sqrt{\frac{\ln t}{N_t(a)}}\right],
\end{equation}
where $c>0$ is an exploration parameter (typically $c=\sqrt{2}$).
\end{definition}

The quantity $Q_t(a)+c\sqrt{(\ln t)/N_t(a)}$ serves as an \emph{upper confidence bound} on the true value $q_*(a)$.  The second term\,---\,the \emph{exploration bonus}\index{exploration bonus}\,---\,is large when action $a$ has been tried infrequently, encouraging the agent to explore it.  The factor $\ln t$ in the numerator grows slowly, ensuring that the agent continues to revisit uncertain arms as the horizon grows, while the $N_t(a)$ in the denominator ensures that the bonus shrinks as an arm is sampled more.

\paragraph{Derivation from Hoeffding's inequality.}
\index{Hoeffding's inequality}
For bounded rewards $R\in[0,1]$, Hoeffding's inequality gives
\[
\Prob\!\big(q_*(a) > Q_t(a) + u\big) \le e^{-2N_t(a)\,u^2}.
\]
Setting the right-hand side equal to $t^{-4}$ (so that the failure probability is summable) and solving for $u$ yields
\[
u = \sqrt{\frac{2\ln t}{N_t(a)}},
\]
which motivates the form~\eqref{eq:ucb1} with $c=\sqrt{2}$.

\subsection*{Worked Numerical Example}

To see UCB1 in action, consider a simple $3$-armed bandit with unknown true values.  We initialise each arm with one forced pull (a common initialisation strategy to avoid division by zero).

\begin{example}[UCB1 arm-selection trace]
\label{ex:ucb1-trace}
\index{UCB1!worked example}
Suppose $k=3$ and after the three forced initial pulls ($t=1,2,3$) we have observed rewards $R_1=0.8$ from arm~1, $R_2=0.4$ from arm~2, and $R_3=0.6$ from arm~3.  Thus $Q_3(1)=0.8$, $Q_3(2)=0.4$, $Q_3(3)=0.6$, and $N_3(a)=1$ for all $a$.  Set $c=\sqrt{2}$.

At $t=4$, the UCB scores are:
\begin{align*}
U_4(1) &= 0.8 + \sqrt{\tfrac{2\ln 4}{1}} \approx 0.8 + 1.664 = 2.464,\\
U_4(2) &= 0.4 + \sqrt{\tfrac{2\ln 4}{1}} \approx 0.4 + 1.664 = 2.064,\\
U_4(3) &= 0.6 + \sqrt{\tfrac{2\ln 4}{1}} \approx 0.6 + 1.664 = 2.264.
\end{align*}
Arm~1 is selected.  Suppose we observe $R_4=0.3$.  Now $Q(1)=(0.8+0.3)/2=0.55$ and $N(1)=2$.

At $t=5$:
\begin{align*}
U_5(1) &= 0.55 + \sqrt{\tfrac{2\ln 5}{2}} \approx 0.55 + 1.270 = 1.820,\\
U_5(2) &= 0.4 + \sqrt{\tfrac{2\ln 5}{1}} \approx 0.4 + 1.794 = 2.194,\\
U_5(3) &= 0.6 + \sqrt{\tfrac{2\ln 5}{1}} \approx 0.6 + 1.794 = 2.394.
\end{align*}
Arm~3 is selected.  The key observation is that arm~1, despite having the highest initial reward, was ``punished'' by the outcome $R_4=0.3$, and arms~2 and~3\,---\,which have not been pulled again\,---\,have maintained their large exploration bonuses.  The algorithm will continue cycling until all arms have been sampled enough that the bonuses shrink and the empirical means dominate.
\end{example}

The figure below illustrates how confidence intervals shrink as an arm is pulled more, with the UCB score converging toward the true value.

\begin{figure}[ht]
\centering
\begin{tikzpicture}[>=Stealth, font=\small]
  % Axes
  \draw[->] (0,0) -- (9.5,0) node[right] {$N_t(a)$};
  \draw[->] (0,0) -- (0,4.5) node[above] {value};

  % True value line
  \draw[dashed, gray, thick] (0.2,2) -- (9.2,2) node[right, gray] {$q_*(a)=0.6$};

  % Empirical mean line (converging to true value)
  \draw[blue!70!black, thick] plot[smooth] coordinates {
    (0.5,3.4)(1.5,2.9)(2.5,2.6)(3.5,2.4)(4.5,2.3)(5.5,2.2)(6.5,2.15)(7.5,2.1)(8.5,2.05)
  } node[right, blue!70!black] {$U_t(a)$};

  % Mean estimate
  \draw[orange!80!black, thick] plot[smooth] coordinates {
    (0.5,0.8)(1.5,1.1)(2.5,1.3)(3.5,1.5)(4.5,1.6)(5.5,1.7)(6.5,1.8)(7.5,1.85)(8.5,1.9)
  };
  \node[orange!80!black] at (8.5,1.7) {$Q_t(a)$};

  % Shaded confidence band
  \fill[blue!10] plot[smooth] coordinates {
    (0.5,0.8)(1.5,1.1)(2.5,1.3)(3.5,1.5)(4.5,1.6)(5.5,1.7)(6.5,1.8)(7.5,1.85)(8.5,1.9)
  } -- plot[smooth] coordinates {
    (8.5,2.05)(7.5,2.1)(6.5,2.15)(5.5,2.2)(4.5,2.3)(3.5,2.4)(2.5,2.6)(1.5,2.9)(0.5,3.4)
  } -- cycle;

  % Vertical brace annotations
  \draw[<->, gray!60] (1.5,1.1) -- (1.5,2.9) node[midway, right=2pt, gray!80] {\tiny wide};
  \draw[<->, gray!60] (7.5,1.85) -- (7.5,2.1) node[midway, right=2pt, gray!80] {\tiny narrow};

  % x-axis ticks
  \foreach \x/\lab in {0.5/1, 1.5/2, 2.5/5, 3.5/10, 4.5/20, 5.5/30, 6.5/50, 7.5/80, 8.5/120}{
    \draw (\x, -0.07) -- (\x, 0.07);
    \node[below] at (\x,-0.12) {\lab};
  }
\end{tikzpicture}
\caption{Schematic showing UCB confidence bounds shrinking as arm $a$ is pulled more often.  The blue shaded band represents the confidence interval $[Q_t(a),\, U_t(a)]$.  As $N_t(a)$ grows, the exploration bonus $c\sqrt{\ln t/N_t(a)}$ decays and the UCB score converges toward the true value $q_*(a)$ (dashed line).}
\label{fig:ucb-shrinking}
\end{figure}

\subsection*{Regret Bound and Proof Sketch}

\begin{theorem}[Regret bound for UCB1]
\label{thm:ucb1-regret}
For a $k$-armed bandit with rewards in $[0,1]$, the UCB1 algorithm with $c=\sqrt{2}$ satisfies
\begin{equation}
\label{eq:ucb1-regret}
\E[\cR_T] \le \sum_{a:\Delta_a>0}\left(\frac{8\ln T}{\Delta_a}\right) + \left(1+\frac{\pi^2}{3}\right)\sum_{a=1}^{k}\Delta_a.
\end{equation}
In particular, $\E[\cR_T]=O\!\left(\sqrt{kT\ln T}\right)$ in the worst case over gaps $\Delta_a$.
\end{theorem}

The proof of Theorem~\ref{thm:ucb1-regret} is instructive because it illustrates the key technique of \emph{counting bad events}.  We sketch the main argument here.

\begin{proof}[Proof sketch]
The regret decomposition~\eqref{eq:regret-decomposition} tells us that $\E[\cR_T]=\sum_{a:\Delta_a>0}\Delta_a\,\E[N_T(a)]$.  It therefore suffices to bound $\E[N_T(a)]$ for each sub-optimal arm $a$.

Fix a sub-optimal arm $a$ with gap $\Delta_a>0$ and let $\ell = \lceil 8\ln T/\Delta_a^2\rceil$.  Arm $a$ is pulled at time $t$ only if its UCB score exceeds that of the optimal arm $a^*$:
\[
Q_t(a) + c\sqrt{\frac{\ln t}{N_t(a)}} \ge Q_t(a^*) + c\sqrt{\frac{\ln t}{N_t(a^*)}}.
\]
For this to happen when $N_t(a)\ge \ell$, either the empirical mean of $a^*$ must be far below its true value, or the empirical mean of $a$ must be far above its true value.  By Hoeffding's inequality (applied with a union bound over all $t$), the probability of either event occurring after $\ell$ pulls is at most $2t^{-4}$, which is summable.  A careful calculation then gives
\[
\E[N_T(a)] \le \ell + \sum_{t=1}^{T}\Prob(\text{bad event at }t) \le \frac{8\ln T}{\Delta_a^2} + 1 + \frac{\pi^2}{3}.
\]
Multiplying by $\Delta_a$ and summing over sub-optimal arms yields~\eqref{eq:ucb1-regret}.
\end{proof}

\begin{remark}
The key insight is that UCB achieves $O(\log T)$ regret\,---\,dramatically better than the $\Omega(T)$ linear regret of $\varepsilon$-greedy with fixed $\varepsilon$.  Moreover, the exploration is directed rather than uniform: arms that are more uncertain or have higher empirical means receive more attention.  The $\pi^2/3$ constant in~\eqref{eq:ucb1-regret} arises as follows: the Hoeffding failure probability per round is $2t^{-4}$, but summing over all possible values of the two sample sizes (each up to~$t$) introduces a factor of~$t^2$, yielding a per-round contribution of $t^2 \cdot 2t^{-4} = 2t^{-2}$.  Summing gives $2\sum_{t=1}^{\infty}t^{-2}=2\cdot\pi^2/6=\pi^2/3$.
\end{remark}


% ===========================================================
\section{Thompson Sampling}
\label{sec:thompson}
\index{Thompson Sampling}

An alternative approach to the exploration--exploitation trade-off is to maintain a \emph{Bayesian posterior}\index{Bayesian posterior} over the action values and use it to guide exploration.  Rather than constructing an explicit optimistic bound, we let the posterior's uncertainty do the work automatically.

\subsection*{The Posterior-Sampling Intuition}

Thompson Sampling\,---\,introduced by William R.\ Thompson in 1933, long before the term ``bandit problem'' existed\,---\,is based on a beautifully simple idea: act as if the world were as good as it might plausibly be.  At each step, we draw a single sample from the posterior distribution of each arm's value and then act greedily with respect to those samples.

Why does this work?  Consider two arms, one well-explored (narrow posterior) and one poorly explored (wide posterior).  For the well-explored arm, the sample will almost surely fall close to the true mean.  For the poorly explored arm, the sample has a non-negligible chance of being much higher than the true mean.  This means uncertain arms are explored with higher probability, naturally implementing the optimism principle\,---\,but in a randomised, Bayesian fashion rather than a worst-case frequentist one.

Crucially, the probability that arm $a$ is selected equals the posterior probability that $a$ is optimal:
\[
\Prob(A_t = a) = \Prob\!\bigl(\tilde{q}(a) = \max_{b}\tilde{q}(b)\bigr) \approx \Prob\!\bigl(q_*(a) = \max_{b}q_*(b)\;\big|\;\text{data}\bigr).
\]
This is the \emph{probability matching}\index{probability matching} property: the agent samples each arm in proportion to how likely it is to be optimal given everything observed so far.  As evidence accumulates, these probabilities concentrate on the true best arm.

\begin{definition}[Thompson Sampling]
\label{def:thompson}
\textbf{Thompson Sampling} (Thompson, 1933) proceeds as follows at each step:
\begin{enumerate}
\item For each arm $a$, sample $\tilde{q}(a)$ from the posterior distribution of $q_*(a)$.
\item Play the arm with the highest sample: $A_t = \argmax_{a}\tilde{q}(a)$.
\item Update the posterior using the observed reward.
\end{enumerate}
\end{definition}

\paragraph{Bernoulli bandits with Beta priors.}
\index{Beta-Bernoulli bandit}
Consider the case where rewards are Bernoulli with unknown success probability $\theta_a$ for each arm.  If we place a $\text{Beta}(\alpha_a, \beta_a)$ prior on $\theta_a$, then after observing $s_a$ successes and $f_a$ failures, the posterior is
\[
\theta_a \mid \text{data} \;\sim\; \text{Beta}(\alpha_a + s_a,\;\beta_a + f_a).
\]
Thompson Sampling then draws $\tilde\theta_a\sim\text{Beta}(\alpha_a+s_a,\beta_a+f_a)$ for each arm and plays $A_t=\argmax_a\tilde\theta_a$.

The Beta distribution is conjugate to the Bernoulli likelihood\,---\,meaning posterior updates have a simple closed form.  This makes the algorithm extremely efficient: no numerical integration is needed, and each update requires only incrementing two counters.

\begin{algorithm}[ht]
\caption{Thompson Sampling for Bernoulli Bandits}
\label{alg:thompson}
\KwIn{Number of arms $k$, prior parameters $\alpha_0,\beta_0$ (e.g.\ $\alpha_0=\beta_0=1$)}
Initialise $S(a)\leftarrow 0$, $F(a)\leftarrow 0$ for all $a$\;
\For{$t=1,2,\dots,T$}{
  \For{$a=1,\dots,k$}{
    Sample $\tilde\theta_a\sim\mathrm{Beta}(\alpha_0+S(a),\;\beta_0+F(a))$\;
  }
  $A_t\leftarrow\argmax_a\,\tilde\theta_a$\;
  Observe reward $R_t\in\{0,1\}$\;
  \If{$R_t=1$}{$S(A_t)\leftarrow S(A_t)+1$}
  \Else{$F(A_t)\leftarrow F(A_t)+1$}
}
\end{algorithm}

\subsection*{Step-by-Step Worked Example}

\begin{example}[Thompson Sampling with three Bernoulli arms]
\label{ex:thompson-worked}
\index{Thompson Sampling!worked example}
Consider a $3$-armed Bernoulli bandit with true success probabilities $\theta_1=0.3$, $\theta_2=0.5$, $\theta_3=0.7$.  We begin with uniform priors $\mathrm{Beta}(1,1)$ on each arm.

\textbf{Round 1.}  We draw one sample from each $\mathrm{Beta}(1,1)=\mathrm{Uniform}(0,1)$ prior.  Suppose the samples are $\tilde\theta_1=0.62$, $\tilde\theta_2=0.29$, $\tilde\theta_3=0.45$.  We select arm~1 (highest sample).  The outcome is a failure ($R=0$), so we update: arm~1 posterior becomes $\mathrm{Beta}(1,2)$.

\textbf{Round 2.}  We draw from $\mathrm{Beta}(1,2)$, $\mathrm{Beta}(1,1)$, $\mathrm{Beta}(1,1)$.  $\mathrm{Beta}(1,2)$ is skewed toward smaller values (mean $1/3$), so arm~1 is now less likely to win.  Suppose $\tilde\theta_1=0.21$, $\tilde\theta_2=0.74$, $\tilde\theta_3=0.38$.  Arm~2 is selected.  Suppose success ($R=1$): arm~2 becomes $\mathrm{Beta}(2,1)$.

\textbf{Round 3.}  Now arm~3 has never been tried and retains its $\mathrm{Beta}(1,1)$ prior.  Arm~2 has a $\mathrm{Beta}(2,1)$ posterior (mean $2/3$, skewed upward).  Suppose the samples are $\tilde\theta_1=0.18$, $\tilde\theta_2=0.55$, $\tilde\theta_3=0.82$, so arm~3 is selected.  Suppose success: arm~3 becomes $\mathrm{Beta}(2,1)$.

After many rounds, arms with higher true probabilities accumulate more successes, their Beta posteriors shift rightward, and their samples are stochastically larger.  The posteriors concentrate, and arm~3 is eventually selected almost exclusively.
\end{example}

Figure~\ref{fig:beta-posteriors} illustrates how the Beta posteriors evolve for each arm as data accumulates.

\begin{figure}[ht]
\centering
\begin{tikzpicture}[>=Stealth, font=\small, scale=0.92]
  % Panel labels and axes for three arms
  % We draw schematic Beta densities at t=1 and t=50

  % --- Arm 1: theta=0.3 ---
  \begin{scope}[xshift=0cm]
    \draw[->] (0,0) -- (3.2,0) node[right] {$\theta_1$};
    \draw[->] (0,0) -- (0,3.2);
    \node[above] at (1.5,3.2) {\small Arm 1 ($\theta_1=0.3$)};
    % Prior Beta(1,1): flat
    \draw[gray, dashed, thick] (0.1,1.8) -- (3.0,1.8) node[right, gray]{\tiny prior};
    % Posterior Beta(~8,18) concentrated near 0.3: draw a rough bell
    \draw[blue!70!black, thick, domain=0.1:3.0, samples=40]
      plot (\x, {2.8*exp(-8*(\x/3-0.3)^2)});
    \node[blue!70!black] at (2.4,0.5) {\tiny posterior};
    \draw (0.9,-0.08) -- (0.9,0.08) node[below=3pt]{\tiny $0.3$};
  \end{scope}

  % --- Arm 2: theta=0.5 ---
  \begin{scope}[xshift=4cm]
    \draw[->] (0,0) -- (3.2,0) node[right] {$\theta_2$};
    \draw[->] (0,0) -- (0,3.2);
    \node[above] at (1.5,3.2) {\small Arm 2 ($\theta_2=0.5$)};
    \draw[gray, dashed, thick] (0.1,1.8) -- (3.0,1.8) node[right, gray]{\tiny prior};
    \draw[blue!70!black, thick, domain=0.1:3.0, samples=40]
      plot (\x, {2.8*exp(-8*(\x/3-0.5)^2)});
    \node[blue!70!black] at (2.4,0.5) {\tiny posterior};
    \draw (1.5,-0.08) -- (1.5,0.08) node[below=3pt]{\tiny $0.5$};
  \end{scope}

  % --- Arm 3: theta=0.7 ---
  \begin{scope}[xshift=8cm]
    \draw[->] (0,0) -- (3.2,0) node[right] {$\theta_3$};
    \draw[->] (0,0) -- (0,3.2);
    \node[above] at (1.5,3.2) {\small Arm 3 ($\theta_3=0.7$)};
    \draw[gray, dashed, thick] (0.1,1.8) -- (3.0,1.8) node[right, gray]{\tiny prior};
    \draw[blue!70!black, thick, domain=0.1:3.0, samples=40]
      plot (\x, {2.8*exp(-8*(\x/3-0.7)^2)});
    \node[blue!70!black] at (0.6,0.5) {\tiny posterior};
    \draw (2.1,-0.08) -- (2.1,0.08) node[below=3pt]{\tiny $0.7$};
  \end{scope}
\end{tikzpicture}
\caption{Schematic evolution of Beta posteriors for a $3$-armed Bernoulli bandit.  Dashed lines show the flat $\mathrm{Beta}(1,1)$ prior; solid curves show the posterior after $T=50$ rounds.  Each posterior concentrates near the true success probability of its arm.  The posterior for arm~3 (the best arm) is the tallest and narrowest, reflecting the most accumulated evidence.}
\label{fig:beta-posteriors}
\end{figure}

\subsection*{Gaussian Bandits Variant}

Thompson Sampling extends naturally to Gaussian reward distributions.  Suppose $R_t \mid A_t=a \sim \cN(\mu_a, \sigma^2)$ with known variance $\sigma^2$ and unknown mean $\mu_a$.  A conjugate prior is $\mu_a \sim \cN(\mu_0, \sigma_0^2)$.  After $n$ observations with sample mean $\bar{r}_a$, the posterior is
\[
\mu_a \mid \text{data} \;\sim\; \cN\!\left(\frac{\sigma^{-2}\cdot n\,\bar{r}_a + \sigma_0^{-2}\mu_0}{\sigma^{-2}\cdot n + \sigma_0^{-2}},\;\frac{1}{\sigma^{-2}\cdot n + \sigma_0^{-2}}\right).
\]
At each step, we sample $\tilde\mu_a$ from this posterior for each arm and select $A_t=\argmax_a\tilde\mu_a$.  The Gaussian case provides a continuous analogue of the Beta-Bernoulli model and is convenient when rewards are unbounded.

\subsection*{Thompson Sampling vs.\ UCB}

Both algorithms achieve $O(\log T)$ regret, but they differ in character.  UCB is a \emph{deterministic} algorithm (given the observations, it always makes the same choice) and its optimism is explicit via the confidence term.  Thompson Sampling is \emph{randomised}: even if the current samples happen not to select the best arm, the probability that this happens decreases as evidence accumulates.

In practice, Thompson Sampling often achieves lower regret than UCB1, particularly in the early phase when priors are informative.  Its main limitation is the need to specify a prior and to sample from the posterior, which can be computationally expensive for complex models.  UCB, by contrast, requires no prior and is straightforward to implement in any setting where empirical averages are meaningful.

\begin{remark}
Thompson Sampling has several attractive properties:
\begin{itemize}[leftmargin=*]
\item It naturally balances exploration and exploitation: uncertain arms are explored because their posterior is wide, producing high samples with non-negligible probability.
\item Empirically, it often matches or outperforms UCB.
\item It achieves $O(\log T)$ Bayesian regret (Kaufmann, Korda, and Munos, 2012) and $O(\sqrt{kT\ln k})$ frequentist regret in many settings.
\item The algorithm is modular: changing the reward model requires only changing the likelihood and the conjugate prior, leaving the high-level loop unchanged.
\end{itemize}
\end{remark}


% ===========================================================
\section{The Lai--Robbins Lower Bound}
\label{sec:lai-robbins}
\index{Lai--Robbins lower bound}

We have seen that UCB1 and Thompson Sampling both achieve $O(\log T)$ regret.  A natural question arises: is $O(\log T)$ the best possible, or could a cleverer algorithm do even better?  The landmark result of Lai and Robbins (1985) answers this question definitively\,---\,no consistent algorithm can do better than $\Omega(\log T)$.

\subsection*{Why $\log T$ Is a Fundamental Limit}

The intuitive reason for the $\log T$ barrier is that any algorithm that eventually learns which arm is best must \emph{continuously test} the hypothesis that a given arm is sub-optimal.  Distinguishing arm $a$ (with mean $\mu_a$) from the optimal arm (with mean $\mu^*>\mu_a$) is a statistical testing problem.  By standard information-theoretic arguments, if we have seen only $n$ samples from arm $a$, we can distinguish $\mu_a$ from $\mu^*$ only if $n \gtrsim 1/\mathrm{KL}(\nu_a\|\nu_{a^*})$, where the KL divergence\index{KL divergence} quantifies how different the two distributions look per sample.  For the algorithm to be \emph{consistent}\,---\,meaning it eventually identifies the best arm on every problem instance\,---\,it must pull arm $a$ at least $\Omega(\log T/\mathrm{KL}(\nu_a\|\nu_{a^*}))$ times.  Multiplying by the gap $\Delta_a$ gives the lower bound~\eqref{eq:lai-robbins} below.

\begin{definition}[Consistent algorithm]
\label{def:consistent}
\index{consistent algorithm}
A bandit algorithm is \textbf{consistent} if, for every problem instance and every sub-optimal arm $a$,
\[
\E[N_T(a)] = o(T^\alpha) \quad \text{for all } \alpha \in (0,1).
\]
In words, the expected number of pulls of any sub-optimal arm grows sub-polynomially in~$T$, so the fraction of time spent on it tends to zero.  Consistency is a minimal rationality requirement: a consistent algorithm eventually stops wasting time on clearly inferior actions.  Note that any algorithm with $O(\log T)$ regret (such as UCB1 or Thompson Sampling) is automatically consistent.
\end{definition}

\begin{theorem}[Lai and Robbins, 1985]
\label{thm:lai-robbins}
For any consistent bandit algorithm and any problem instance with sub-optimality gaps $\Delta_a>0$,
\begin{equation}
\label{eq:lai-robbins}
\liminf_{T\to\infty} \frac{\E[\cR_T]}{\ln T} \ge \sum_{a:\Delta_a>0} \frac{\Delta_a}{\mathrm{KL}(\nu_a\|\nu_{a^*})},
\end{equation}
where $\mathrm{KL}(\nu_a\|\nu_{a^*})$ denotes the Kullback--Leibler divergence between the reward distributions of arm $a$ and the optimal arm $a^*$.
\end{theorem}

\subsection*{Information-Theoretic Argument Sketch}

The proof of Theorem~\ref{thm:lai-robbins} can be understood through a \emph{change-of-measure} argument.  Suppose we have a bandit problem $\cP$ and we consider an alternative problem $\cP'$ that differs only in the reward distribution of arm $a$: in $\cP'$, arm $a$ is actually the best arm.  Any consistent algorithm must behave differently on $\cP$ and $\cP'$, pulling arm $a$ less often in $\cP$ (where it is sub-optimal) than in $\cP'$ (where it is optimal).  The divergence $\mathrm{KL}(\nu_a\|\nu_{a^*})$ controls how quickly the algorithm can distinguish $\cP$ from $\cP'$ given samples from arm $a$.  A data-processing inequality then establishes that the number of pulls of arm $a$ in $\cP$ must grow at least as $\ln T / \mathrm{KL}(\nu_a\|\nu_{a^*})$ for the algorithm to reliably distinguish the two problems.

\subsection*{Worked Example: Gaussian Case}

For Gaussian rewards with unit variance, $\mathrm{KL}(\nu_a\|\nu_{a^*})=\Delta_a^2/2$, and the bound becomes
\[
\liminf_{T\to\infty}\frac{\E[\cR_T]}{\ln T}\ge \sum_{a:\Delta_a>0}\frac{2}{\Delta_a}.
\]

\begin{example}[Gaussian lower bound, two arms]
\label{ex:gaussian-lower-bound}
\index{Lai--Robbins lower bound!Gaussian case}
Consider a $2$-armed Gaussian bandit with $q_*(1)=0$ and $q_*(2)=\Delta>0$ (unit variances).  The lower bound gives
\[
\liminf_{T\to\infty}\frac{\E[\cR_T]}{\ln T}\ge \frac{2}{\Delta}.
\]
Notice that when $\Delta$ is small\,---\,the two arms are hard to distinguish\,---\,the bound is large: the algorithm must pull the sub-optimal arm many times to be sure it is not actually the best.  Conversely, when $\Delta$ is large, the arms are easy to tell apart and only $O(1/\Delta)$ pulls are needed, incurring proportionally smaller regret per forced exploration.  This formalises the intuition that bandit problems with nearly-equal arm values are intrinsically harder than those with clearly separated values.
\end{example}

\begin{remark}
The Lai--Robbins bound implies that $\Omega(\log T)$ regret is unavoidable for any algorithm that eventually learns the optimal action.  Both UCB and Thompson Sampling achieve regret of order $O(\log T)$, so they are \emph{asymptotically optimal}\index{asymptotically optimal} in the instance-dependent sense.  The KL-UCB algorithm (Garivier and Capp\'{e}, 2011) and Thompson Sampling with conjugate priors are known to achieve the exact Lai--Robbins constant, not just the correct rate.
\end{remark}


% ===========================================================
\section{Contextual Bandits}
\label{sec:contextual}
\index{contextual bandit}

The bandit formulation studied so far is \emph{context-free}: the agent observes only the rewards, not any side information about the environment.  In most real-world applications, however, rich context is available before each decision.  A recommendation system knows the user's browsing history; a clinical trial knows the patient's age and biomarkers; an ad-placement system knows the content of the webpage being viewed.  Ignoring this context leads to a one-size-fits-all policy, whereas incorporating it enables personalisation.  This motivates the \emph{contextual bandit} problem.

\begin{definition}[Contextual bandit]
\label{def:contextual-bandit}
At each step $t$:
\begin{enumerate}
\item A context $x_t\in\mathcal{X}$ is revealed to the agent.
\item The agent selects $A_t\in\cA$.
\item A reward $R_t$ is observed, drawn from a distribution that depends on both $x_t$ and $A_t$.
\end{enumerate}
The goal is to learn a policy $\pi:\mathcal{X}\to\cA$ that maximises the expected cumulative reward.
\end{definition}

The key difference from the standard bandit: the optimal action may vary with the context.  In a drug trial, a treatment that is best for elderly patients may be harmful for young ones.  A good contextual-bandit algorithm must learn this dependence from data.

\subsection*{LinUCB: Upper Confidence Bounds with Linear Models}

A widely used contextual-bandit algorithm is \textbf{LinUCB}\index{LinUCB} (Li et al., 2010), which assumes that the expected reward is linear in the context:
\begin{equation}
\label{eq:linucb-model}
\E[R_t \mid x_t, A_t=a] = x_t^\top \theta_a^*,
\end{equation}
where $\theta_a^*\in\R^d$ is an unknown parameter vector for arm $a$.  Given observations, the least-squares estimate $\hat\theta_a$ can be computed, and ridge regression provides a confidence ellipsoid around it.

\begin{definition}[LinUCB decision rule]
\label{def:linucb}
At each step $t$, LinUCB selects
\begin{equation}
\label{eq:linucb}
A_t = \argmax_{a\in\cA}\left[x_t^\top\hat\theta_a(t) + \alpha\,\sqrt{x_t^\top A_a(t)^{-1}x_t}\right],
\end{equation}
where $A_a(t) = \lambda I + \sum_{s<t:\,A_s=a}x_s x_s^\top$ is the regularised design matrix for arm $a$ and $\alpha>0$ is an exploration parameter.
\end{definition}

The term $\sqrt{x_t^\top A_a(t)^{-1}x_t}$ is the \emph{ellipsoidal exploration bonus}\index{ellipsoidal exploration bonus}: it is large when the context $x_t$ lies in a direction that has not been well-explored for arm $a$, and it shrinks as more diverse contexts are observed.  LinUCB enjoys a regret bound of $O(d\sqrt{T\ln T})$ under the linearity assumption~\eqref{eq:linucb-model}, where $d$ is the context dimension.

\subsection*{Practical Applications}

Contextual bandits arise naturally in a number of high-stakes domains.

\paragraph{Online advertising.}
\index{online advertising}
When a user visits a webpage, the system must decide which advertisement to display.  The context $x_t$ includes user features (demographics, browsing history) and page features (topic, time of day).  Each arm corresponds to an ad, and the reward is a click indicator.  Because user preferences vary enormously, a personalised policy can significantly outperform one that displays the same ad to everyone.  LinUCB was deployed by Yahoo!\ for news article recommendation (Li et al., 2010), yielding substantial click-through-rate improvements over non-contextual approaches.

\paragraph{Clinical trials and personalised medicine.}
\index{clinical trials}
\index{personalised medicine}
In adaptive clinical trials, a patient arrives with a covariate vector $x_t$ (age, weight, genetic markers) and the clinician must assign a treatment from a set $\cA$.  The outcome $R_t$ is a health measure such as reduction in tumour size or blood pressure.  A contextual bandit policy learns, over the course of the trial, which treatment is best for which patient subgroup.  Compared to a classical randomised controlled trial (which assigns treatments without using patient covariates), a contextual bandit can reduce the number of patients assigned to inferior treatments while still generating enough data to reach valid conclusions.

\paragraph{News recommendation.}
\index{news recommendation}
A news website must decide which article to show a user on the front page.  The context encodes user history and article features; the reward is reading time or engagement.  The exploration--exploitation trade-off is especially delicate here: showing only articles the user already prefers sacrifices the chance of discovering a new topic they might love.

\subsection*{From Contextual Bandits to Full RL}

Contextual bandits occupy a natural middle ground between the stateless bandit problem and full reinforcement learning.

\begin{description}[leftmargin=!, labelwidth=4cm]
\item[Stateless bandits] No context; one global optimal action.
\item[Contextual bandits] Context $x_t$ observed; optimal action depends on $x_t$; but actions do not influence future contexts.
\item[Full RL (MDP)] State $s_t$ observed; optimal action depends on $s_t$; and actions influence the distribution of future states.
\end{description}

The critical distinction is that in the contextual bandit setting, actions do not affect the distribution of future contexts\,---\,pulling an arm now has no bearing on what context appears next.  In full RL (studied from Chapter~\ref{ch:mdp} onward), actions shape the future, requiring the agent to reason about long-horizon consequences.  Techniques developed for contextual bandits\,---\,especially the UCB-style exploration bonus and posterior-sampling approaches\,---\,have inspired algorithms for RL exploration such as UCRL (Jaksch, Ortner, and Auer, 2010) and RLSVI (Osband et al., 2016).


% ===========================================================
\section{Empirical Comparison}
\label{sec:bandit-experiments}

To ground the theoretical discussion, consider the standard \emph{10-armed Gaussian testbed} introduced by Sutton and Barto:

\begin{example}[Gaussian $10$-armed testbed]
\label{ex:testbed}
\index{10-armed testbed}
Each arm $a=1,\dots,10$ has a true value $q_*(a)\sim\cN(0,1)$, sampled once and then fixed.  Rewards are generated as $R_t\mid A_t=a \sim \cN(q_*(a),1)$.
\end{example}

Results are averaged over many independently sampled problem instances to reduce the influence of any particular configuration of arm values.

\subsection*{Regret Curves}

Figure~\ref{fig:regret-comparison} shows schematic cumulative regret curves for the algorithms studied in this chapter.  Several qualitative features are worth noting.

\begin{figure}[ht]
\centering
\begin{tikzpicture}[>=Stealth, font=\small]
  % Axes
  \draw[->] (0,0) -- (10.5,0) node[right] {$t$};
  \draw[->] (0,0) -- (0,5.0) node[above] {$\E[\cR_t]$};

  % x-axis labels
  \foreach \x/\lab in {0/0, 2.5/250, 5/500, 7.5/750, 10/1000}{
    \draw (\x,-0.08) -- (\x,0.08);
    \node[below] at (\x,-0.12) {\lab};
  }
  % y-axis labels
  \foreach \y/\lab in {0/0, 1/50, 2/100, 3/150, 4/200}{
    \draw (-0.08,\y) -- (0.08,\y);
    \node[left] at (-0.12,\y) {\lab};
  }

  % Pure greedy: linear (steep)
  \draw[red!80!black, very thick] (0,0) -- (10,4.6)
    node[right, red!80!black] {\small Greedy};

  % eps-greedy eps=0.1: roughly linear but shallower
  \draw[orange!80!black, very thick] plot[smooth] coordinates {
    (0,0)(1,0.7)(2,1.2)(3,1.6)(4,2.0)(5,2.35)(6,2.68)(7,2.98)(8,3.24)(9,3.48)(10,3.70)
  } node[right, orange!80!black] {\small $\varepsilon$-greedy ($\varepsilon=0.1$)};

  % UCB1: grows like log T (concave)
  \draw[blue!80!black, very thick] plot[smooth] coordinates {
    (0,0)(0.5,0.4)(1,0.65)(2,0.95)(3,1.15)(4,1.30)(5,1.43)(6,1.54)(7,1.64)(8,1.72)(9,1.80)(10,1.87)
  } node[right, blue!80!black] {\small UCB1};

  % Thompson Sampling: slightly below UCB1
  \draw[green!60!black, very thick] plot[smooth] coordinates {
    (0,0)(0.5,0.35)(1,0.55)(2,0.80)(3,0.97)(4,1.10)(5,1.20)(6,1.29)(7,1.37)(8,1.44)(9,1.50)(10,1.55)
  } node[right, green!60!black] {\small Thompson};

  % Lai-Robbins lower bound (dashed)
  \draw[black!50, thick, dashed] plot[smooth] coordinates {
    (0,0)(1,0.38)(2,0.58)(3,0.71)(4,0.81)(5,0.89)(6,0.96)(7,1.02)(8,1.07)(9,1.12)(10,1.16)
  } node[right, black!50] {\small LR bound};

\end{tikzpicture}
\caption{Schematic cumulative regret $\E[\cR_t]$ versus time $t$ for the algorithms studied in this chapter on the $10$-armed Gaussian testbed.  Pure greedy exhibits linear regret, while $\varepsilon$-greedy achieves lower but still linear regret.  UCB1 and Thompson Sampling achieve logarithmic regret (concave curves), with Thompson Sampling tracking slightly closer to the Lai--Robbins (LR) lower bound.  Actual performance varies with problem instance and parameter settings.}
\label{fig:regret-comparison}
\end{figure}

\paragraph{Pure greedy.} Because pure greedy locks onto an arm early, it accrues regret linearly.  The slope is steep: the algorithm has a substantial probability of being stuck on a sub-optimal arm for the entire horizon.

\paragraph{$\varepsilon$-greedy.} With $\varepsilon=0.1$, ten percent of all steps are purely random, giving the algorithm ample opportunity to discover better arms.  However, this randomness continues indefinitely, so the regret still grows linearly at rate $\varepsilon\,\bar\Delta/k$ (where $\bar\Delta$ is the average sub-optimality gap), just with a smaller coefficient than pure greedy.

\paragraph{UCB1.} The logarithmic regret is clearly visible as the curve bending toward horizontal.  The initial transient is somewhat expensive because UCB must try each arm at least once, and the early exploration bonuses are all large.  Once every arm has been sampled, the algorithm settles into a nearly-greedy regime, exploring only when genuine uncertainty remains.

\paragraph{Thompson Sampling.} Thompson Sampling typically tracks slightly below UCB1 in terms of cumulative regret, sometimes approaching the Lai--Robbins lower bound closely.  This advantage is most pronounced when the prior is well-specified; misspecified priors can degrade performance.

\subsection*{Parameter Sensitivity}

Both UCB1 and $\varepsilon$-greedy expose a parameter ($c$ and $\varepsilon$, respectively) that controls the exploration rate.

\begin{itemize}[leftmargin=*]
\item For UCB1, setting $c$ too small causes under-exploration: sub-optimal arms are abandoned before their values are reliably estimated.  Setting $c$ too large causes over-exploration: the agent spends too much time on unpromising arms.  The theoretically motivated choice $c=\sqrt{2}$ works well across many benchmarks, but practitioners sometimes tune $c$ on held-out data.
\item For $\varepsilon$-greedy, a common heuristic is to decay $\varepsilon_t = \min(1, \varepsilon_0/t)$, which achieves sub-linear (though not logarithmic) regret.  Even with decay, the optimal schedule depends on problem-specific constants that are unknown in advance.
\item Thompson Sampling is comparatively robust to prior misspecification when the prior is uninformative (e.g., $\mathrm{Beta}(1,1)$), and it does not expose an explicit exploration parameter.  This makes it more ``plug-and-play'' in practice.
\end{itemize}

Table~\ref{tab:comparison} summarises the key properties of all algorithms discussed in this chapter.

\begin{table}[ht]
\centering
\caption{Comparison of multi-armed bandit algorithms.}
\label{tab:comparison}
\begin{tabular}{lllll}
\toprule
Algorithm & Regret & Parameter & Prior & Notes \\
\midrule
Pure greedy      & $\Theta(T)$         & ---          & No  & Can lock onto sub-optimal arm \\
$\varepsilon$-greedy (fixed) & $\Theta(\varepsilon T)$ & $\varepsilon$ & No & Simple; linear regret \\
$\varepsilon$-greedy (decay) & $O(\sqrt{T})$  & schedule & No & Suboptimal; schedule-dependent \\
UCB1             & $O(\log T)$         & $c$          & No  & Deterministic; worst-case optimal \\
Thompson Sampling & $O(\log T)$        & ---          & Yes & Randomised; often best empirically \\
\bottomrule
\end{tabular}
\end{table}


% ===========================================================
\section{Summary}
\label{sec:bandit-summary}

The multi-armed bandit is one of the most thoroughly studied problems in sequential decision-making, and for good reason: it is simple enough to be analysed rigorously, yet rich enough to capture the core tension between exploration and exploitation that pervades all of reinforcement learning.

We began with the basic formulation in Section~\ref{sec:bandit-formulation} and established the regret decomposition $\E[\cR_T]=\sum_a\Delta_a\,\E[N_T(a)]$ in Section~\ref{sec:action-values-regret}.  This identity reframes the problem: controlling regret means controlling how often sub-optimal arms are pulled, which requires the algorithm to learn the arm ranking from noisy observations.

The $\varepsilon$-greedy family (Section~\ref{sec:eps-greedy}) provides a simple baseline.  By reserving a fixed fraction of steps for random exploration, the algorithm prevents locking onto a bad arm but pays the price of perpetual random exploration, yielding linear regret.  The incremental update rule (Section~\ref{sec:incremental}) shows that maintaining accurate value estimates can be done in constant time per step, an important practical consideration.

Upper Confidence Bounds (Section~\ref{sec:ucb}) offer a principled solution.  By adding an exploration bonus proportional to $\sqrt{\ln t / N_t(a)}$, UCB1 directs exploration toward arms that are either genuinely good or insufficiently sampled.  The algorithm achieves $O(\log T)$ regret, which Theorem~\ref{thm:ucb1-regret} makes precise through an instance-dependent bound.  The regret proof illustrates a widely applicable technique: decompose bad events by arm and time, bound each event's probability via concentration inequalities, then sum the pieces.

Thompson Sampling (Section~\ref{sec:thompson}) takes a Bayesian perspective, maintaining a posterior distribution over each arm's value and sampling from it to make decisions.  In the Beta-Bernoulli model, posterior updates are trivially cheap (increment two counters), and the probability-matching property ensures that exploration is calibrated to the actual posterior uncertainty.  Thompson Sampling matches UCB1's $O(\log T)$ regret guarantee while often outperforming it empirically, and it extends cleanly to Gaussian bandits and other exponential-family reward models.

The Lai--Robbins lower bound (Section~\ref{sec:lai-robbins}) closes the circle.  No consistent algorithm can achieve regret growing slower than $c\,\ln T$ for a problem-instance-specific constant $c$ determined by the KL divergences between arm distributions.  Both UCB and Thompson Sampling achieve this rate, making them \emph{asymptotically optimal}.  The proof idea\,---\,that distinguishing a sub-optimal arm from the best requires a minimum number of observations dictated by information theory\,---\,is a prototype for lower-bound arguments throughout reinforcement learning theory.

Contextual bandits (Section~\ref{sec:contextual}) show how the bandit framework scales to real-world complexity.  By conditioning decisions on an observed context, algorithms like LinUCB achieve personalised policies that outperform any context-free approach, at the cost of a higher-dimensional parameter estimation problem.  Applications span online advertising, clinical trials, and content recommendation, and the key ideas\,---\,optimism under uncertainty, posterior sampling\,---\,carry over directly from the non-contextual setting.

Looking forward, the bandit problem is not just a simplified warm-up; it is a building block.  Every RL algorithm must, at its core, solve a local version of the exploration--exploitation problem at each state.  The tools developed here\,---\,confidence bounds, Bayesian posteriors, regret analysis, lower bounds via information theory\,---\,will reappear in the analysis of temporal-difference learning, policy gradient methods, and modern deep RL exploration strategies in later chapters.


% ===========================================================
\section*{Exercises}
\addcontentsline{toc}{section}{Exercises}

\begin{exercise}
Prove that an $\varepsilon$-greedy strategy with fixed $\varepsilon>0$ and at least one strictly sub-optimal action satisfies $\E[\cR_T]=\Omega(T)$.
\end{exercise}

\begin{exercise}
Suppose the optimal action has the highest current estimate.  Find:
\begin{enumerate}[label=(\alph*)]
\item the probability of selecting the optimal action under $\varepsilon$-greedy;
\item the probability of selecting any particular sub-optimal action;
\item the expected number of exploratory steps in $T$ rounds.
\end{enumerate}
\end{exercise}

\begin{exercise}
Let $R_1,\dots,R_n$ be i.i.d.\ rewards with mean $q_*(a)$ and variance $\sigma_a^2$.  Prove that the sample average $Q_n(a)=\frac{1}{n}\sum_{i=1}^{n}R_i$ is unbiased and has variance $\sigma_a^2/n$.
\end{exercise}

\begin{exercise}
Derive the explicit form of the constant-step-size update:
\[
Q_{n+1} = (1-\alpha)^n Q_1 + \sum_{i=1}^{n}\alpha(1-\alpha)^{n-i}R_i.
\]
Discuss how the weights on past observations behave as $n$ grows.
\end{exercise}

\begin{exercise}
Construct a specific two-armed bandit instance in which the pure greedy algorithm, starting from $Q_1(1)=Q_1(2)=0$, locks onto the sub-optimal arm forever with positive probability.  Specify the reward distributions explicitly.
\end{exercise}

\begin{exercise}
Show that any algorithm satisfying $\Prob(A_t=a)\ge\delta$ for all $a$, all $t$, and some $\delta>0$ must have $\E[\cR_T]\ge cT$ for some constant $c>0$ (assuming at least one arm is strictly sub-optimal).
\end{exercise}

\begin{exercise}[UCB exploration bonus]
\label{ex:ucb-bonus}
Starting from Hoeffding's inequality for bounded random variables, derive the UCB1 exploration bonus $\sqrt{2\ln t/N_t(a)}$.  \emph{Hint}: set the failure probability to $t^{-4}$ and take a union bound over arms.
\end{exercise}

\begin{exercise}[Thompson Sampling simulation]
\label{ex:thompson-sim}
Consider a $3$-armed Bernoulli bandit with success probabilities $\theta_1=0.3$, $\theta_2=0.5$, $\theta_3=0.7$.  Using Beta$(1,1)$ priors, simulate Thompson Sampling for $T=1000$ steps.  Plot the posterior distributions of $\theta_1,\theta_2,\theta_3$ at $t=10$, $t=100$, and $t=1000$.  Observe how the posteriors concentrate around the true values.
\end{exercise}

\begin{exercise}
\label{ex:lai-robbins-gaussian}
For a Gaussian bandit with unit variance, the KL divergence is $\mathrm{KL}(\cN(\mu_a,1)\|\cN(\mu^*,1))=\Delta_a^2/2$.  Substitute into the Lai--Robbins bound~\eqref{eq:lai-robbins} and interpret: how does the difficulty of distinguishing arm~$a$ from the optimum depend on~$\Delta_a$?
\end{exercise}

\begin{exercise}[LinUCB derivation]
\label{ex:linucb}
Assume rewards satisfy $R_t = x_t^\top\theta^* + \eta_t$ where $\eta_t\sim\cN(0,\sigma^2)$ and all arms share a single parameter $\theta^*\in\R^d$.  Show that the ridge regression estimate $\hat\theta_t = A_t^{-1}b_t$ (where $A_t=\lambda I + \sum_{s\le t}x_{A_s,s}x_{A_s,s}^\top$ and $b_t=\sum_{s\le t}R_s x_{A_s,s}$) is the maximum-a-posteriori estimate under a Gaussian prior $\theta^*\sim\cN(0,\sigma^2\lambda^{-1}I)$.  Verify that the exploration bonus $\sqrt{x_t^\top A_t^{-1}x_t}$ is the posterior standard deviation of $x_t^\top\theta^*$.
\end{exercise}
